\chapter{ENT}

\medterm{Acoustic Neuroma}-- Acoustic neuroma, also known as vestibular schwannoma, is a benign, slow-growing tumor arising from the Schwann cells of the vestibular portion of the eighth cranial nerve (vestibulocochlear nerve), typically within the internal auditory canal or the cerebellopontine angle, and it accounts for approximately eight to ten percent of all intracranial tumors. The tumor usually presents with unilateral sensorineural hearing loss, tinnitus, and disequilibrium, and less commonly with vertigo, facial numbness (from trigeminal nerve compression), or facial weakness (from facial nerve compression).

\medterm{Adenoidectomy}-- The surgical removal of the pharyngeal tonsils (adenoids), which are lymphoid tissue structures located in the nasopharynx, and it is commonly performed in children for obstructive sleep apnea, chronic nasal obstruction, recurrent otitis media, and chronic rhinosinusitis that is refractory to medical management. The adenoids are typically most prominent between the ages of three and seven years and undergo spontaneous involution during childhood, and adenoid hypertrophy can

\medterm{Allergic Rhinitis}-- Allergic rhinitis, commonly known as hay fever, is a type I IgE-mediated hypersensitivity reaction of the nasal mucosa to inhaled allergens, characterized by sneezing, nasal congestion, rhinorrhea (clear, watery nasal discharge), nasal pruritus, and conjunctival symptoms including itching, tearing, and redness.

\medterm{Anosmia}-- The complete loss of the sense of smell, while hyposmia: a reduced sense of smell, and these conditions can significantly impair quality of life by reducing the ability to appreciate food flavors, detect environmental hazards such as gas leaks and smoke, and enjoy activities. The causes of anosmia are numerous and include post-viral olfactory loss (the most common cause, often following upper respiratory viral infections including SARSCoV-2), chronic rhinosinusitis with nasal polyps, head trauma (particularly frontal lobe injuries that disrupt the olfactory nerves), neurodegenerative diseases (Alzheimer disease, Parkinson disease), congenital conditions (Kallmann syndrome), endocrine disorders (hypothyroidism, diabetes), medication side effects, and toxic exposures (smoking, industrial chemicals). The diagnosis involves olfactory testing using standardized smell identification and threshold tests, nasal endoscopy to evaluate for anatomical causes, and MRI to assess the olfactory bulbs and frontal lobes. Treatment depends on the underlying cause and may include intranasal corticosteroids for inflammatory causes, olfactory training for postviral anosmia, and management of underlying medical conditions.

\medterm{Aphthous Ulcer}-- Recurrent, painful, round or oval ulcers of the oral mucosa, commonly called canker sores. Three types: minor (most common, <1 cm, heal in 7-14 days without scarring), major (>1 cm, deeper, heal over weeks with scarring), and herpetiform (multiple pinhead-sized ulcers that fuse). Etiology unclear: immune-mediated, genetic predisposition, stress, trauma, nutritional deficiencies (B12, folate, iron), and food sensitivities. Diagnosis: clinical. Treatment: symptomatic (topical corticosteroids, amlexanox, lidocaine, oral hygiene). Severe cases: systemic corticosteroids, colchicine, or dapsone.

\medterm{Arytenoid Cartilage}-- The arytenoid cartilages are a pair of small, pyramid-shaped cartilages located on the posterior aspect of the larynx, articulating with the superior border of the cricoid cartilage lamina, and they serve as the attachment points for the vocal cords and the intrinsic laryngeal muscles that control vocal cord movement. Each arytenoid cartilage has a vocal process (anteriorly), to which the vocal cord ligament attaches, and a muscular process (laterally), to which the posterior cricoarytenoid and lateral cricoarytenoid muscles attach. The movement of the arytenoid cartilages on the cricoid cartilage during phonation, breathing, and swallowing causes the vocal cords to adduct (close), abduct (open), and tense, and the posterior cricoarytenoid muscle is the sole abductor of the vocal cords, making it critical for maintaining airway patency. The arytenoid cartilages may become displaced or subluxed by trauma, intubation injury, or surgery, and arytenoid dislocation can cause hoarseness and vocal cord paresis.

\medterm{Audiogram}-- A chart that shows a person’s ability to hear at different pitches or frequencies.

\medterm{Audiologist}-- A health professional who assesses hearing and fits hearing aids.

\medterm{Audiometry}-- A complete hearing test that involves listening to sounds of different frequencies and volume.

\medterm{Cerumen}-- A substance that helps keep dirt out of the ear and lubricates the skin in the ear. More commonly known as earwax.

\medterm{Cerumen Impaction}-- The accumulation of cerumen (earwax) in the external auditory canal to the point where it causes symptoms or prevents adequate visualization of the tympanic membrane, and it is a common condition affecting approximately five to ten percent of children and up to thirty percent of elderly individuals and cognitively impaired adults. Cerumen is a normal, protective secretion of the ceruminous glands in the cartilaginous portion of the external auditory canal, composed of secretions from modified apocrine glands mixed with dead skin cells and hair, and it serves to moisturize the canal skin, trap debris, and provide antimicrobial protection. Symptoms of impaction include aural fullness, decreased hearing, tinnitus, pruritus, pain, and a feeling of obstruction, and impacted cerumen may be visible as a yellow, brown, or dark mass obscuring the tympanic membrane. Treatment includes softening agents such as mineral oil, docusate sodium, or hydrogen peroxide, followed by manual removal using curettes, irrigation, or suction, and patients should be advised against inserting cotton swabs or other objects into the ear canal.

\medterm{Cholesteatoma}-- A pathological accumulation of keratinizing squamous epithelium within the middle ear, mastoid, or temporal bone that behaves like a locally invasive, destructive lesion, causing progressive conductive hearing loss, chronic otorrhea, and potential complications including ossicular erosion, labyrinthine fistula, facial nerve paralysis, and intracranial infections. The condition may be congenital (arising from embryonic epithelial remnants within the temporal bone) or acquired (typically developing from retraction pockets of the pars flaccida or perforations of the tympanic membrane), and it grows by the accumulation and desquamation of keratin debris within the squamous epithelial-lined sac. The cholesteatoma matrix produces proteolytic enzymes and inflammatory mediators that erode adjacent bony structures, including the ossicles, the bony labyrinth, and the mastoid cortex. Diagnosis is confirmed by otoscopy showing a retraction pocket or pearly white mass behind the tympanic

\medterm{Chronic Rhinosinusitis}-- A heterogeneous group of disorders characterized by persistent inflammation of the paranasal sinuses lasting more than twelve weeks, with symptoms including nasal obstruction, purulent rhinorrhea (anterior or posterior), facial pain or pressure, and olfactory dysfunction, and it is classified into chronic rhinosinusitis without nasal polyps and chronic rhinosinusitis with nasal polyps based on the presence or absence of nasal polyps on endoscopic examination. The pathogenesis is multifactorial, involving bacterial biofilm formation, mucociliary dysfunction, immunological dysregulation (type 2 inflammation in polyp disease), and anatomical factors that impair sinus drainage. The diagnosis is confirmed by nasal endoscopy showing mucosal edema, polyps, or purulent drainage from the sinus ostia, and CT imaging of the sinuses showing mucosal thickening and sinus opacification. Treatment includes intranasal corticosteroids, saline irrigation, antibiotics for acute exacerbations, and surgical intervention with functional endoscopic sinus surgery for patients who fail medical therapy.

\medterm{Cochlear Implant}-- A cochlear implant is an electronic medical device that bypasses damaged or absent inner ear hair cells to provide direct electrical stimulation of the auditory nerve, restoring a perception of sound in individuals with severe to profound sensorineural hearing loss who derive limited benefit from hearing aids. The device consists of an external component (microphone, speech processor, and transmitter coil) and an internal component (receiver-stimulator and electrode array inserted into the cochlea), and it works by converting sound into electrical signals that are transmitted to the electrode array, which stimulates the spiral ganglion neurons of the auditory nerve. The implant does not restore normal hearing but provides a representation of environmental sounds and speech that the brain learns to interpret through auditory rehabilitation. Indications include bilateral severe to profound sensorineural hearing loss in adults who receive limited benefit from appropriately fitted hearing aids, and bilateral severe to profound sensorineural hearing loss in children (typically implanted at twelve to eighteen months of age) to support speech and language development. The surgery is performed under general anesthesia and is generally safe, with risks including facial nerve injury, meningitis, and device failure.

\medterm{Conductive Hearing Loss}-- A type of hearing loss caused by impairment of the transmission of sound through the outer ear, middle ear, or both, preventing sound vibrations from reaching the cochlea efficiently, and it is typically characterized by a conductive gap (an air-bone gap on audiometry indicating a difference between air conduction and bone conduction thresholds). Common causes in-

\medterm{Cricoid Cartilage}-- The cricoid cartilage is a complete, ring-shaped cartilage located at the base of the larynx, directly above the trachea and below the thyroid cartilage, and it is the only complete cartilaginous ring in the airway, providing structural support and maintaining airway patency. The cricoid cartilage is shaped like a signet ring, with a narrow anterior arch and a broad posterior lamina that articulates with the arytenoid cartilages superiorly. The cricothyroid membrane, connecting the thyroid and cricoid cartilages anteriorly, is the site of emergency cricothyrotomy (cricothyroidotomy) for securing the airway when intubation is not possible. The cricothyroid muscle, which attaches to the cricoid and thyroid cartilages, is the only intrinsic laryngeal muscle that is not innervated by the recurrent laryngeal nerve, and it functions to tense the vocal cords. The cricoid cartilage is prone to fracture in blunt neck trauma, and cricoid fracture can cause airway compromise, subglottic stenosis, and recurrent laryngeal nerve injury.

\medterm{CSF Rhinorrhea}-- Cerebrospinal fluid rhinorrhea is the leakage of cerebrospinal fluid from the subarachnoid space through a defect in the skull base into the nasal cavity or paranasal sinuses, resulting in a clear, watery nasal discharge that may be mistaken for allergic rhinitis or sinusitis. The most common causes include trauma (skull base fractures), iatrogenic injury (following sinus or skull base surgery), spontaneous

\medterm{Deafness}-- Complete or partial loss of the ability to hear, classified by type: conductive deafness (impaired sound transmission through the external or middle ear, causes include cerumen impaction, otitis media, otosclerosis, tympanic membrane perforation), sensorineural deafness (damage to the cochlea or auditory nerve, causes include presbycusis, noise exposure, ototoxic drugs, Meniere disease, acoustic neuroma), and mixed hearing loss. Diagnosis: audiometry (pure-tone, speech), tympanometry, and otoacoustic emissions. Management: hearing aids, cochlear implants (for severe-to-profound sensorineural loss), bone-anchored hearing aids (for conductive or single-sided deafness), and surgical correction of underlying causes.

\medterm{Decibel}-- A logarithmic unit (dB) used to measure the intensity of sound. The decibel scale compares a given sound pressure level to a reference level (0 dB = 20 micropascals, the threshold of normal hearing at 1 kHz). The scale is logarithmic: an increase of 10 dB represents a tenfold increase in sound intensity. Normal conversation: ~60 dB. Threshold of pain: ~120 dB. Prolonged exposure to sounds >85 dB causes noise-induced hearing loss. Audiometry measures hearing thresholds in dB hearing level (dB HL) across frequencies 250--8000 Hz.

\medterm{Dix-Hallpike Maneuver}-- The Dix-Hallpike maneuver is a diagnostic test used to identify benign paroxysmal positional vertigo by detecting the characteristic nystagmus and vertigo that occur when displaced otoconia within the posterior semicircular canal move in response to gravity during specific head positioning. The patient is seated on the examination table, and the examiner rotates the patient’s head forty-five degrees toward the side being tested, then quickly lays the patient back with the head hanging below the level of the shoulders, maintaining the head rotation.

\medterm{Ear canal}-- A tube leading from the eardrum to the outer ear.

\medterm{Ear Diseases}-- Any disorder affecting the structures of the ear--external ear (pinna, ear canal), middle ear (tympanic membrane, ossicles, Eustachian tube), or inner ear (cochlea, vestibular apparatus). Broad categories: congenital (microtia, atresia, preauricular sinus), infectious (otitis externa, otitis media, mastoiditis), inflammatory (Eustachian tube dysfunction, cholesteatoma), traumatic (tympanic membrane perforation, temporal bone fracture, barotrauma), neoplastic (acoustic neuroma, glomus tumor, squamous cell carcinoma), and degenerative (presbycusis, Meniere disease, otosclerosis). Diagnosis: otoscopy, audiometry, tympanometry, and imaging (CT temporal bone, MRI internal auditory meatus). Management is etiology-specific.

\medterm{Ear Infections}-- Infections of the ear, most commonly affecting the middle ear (otitis media) or the outer ear canal (otitis externa). They are often caused by bacteria or viruses and present with ear pain, drainage, hearing loss, and sometimes fever. Recurrent ear infections in children may require tympanostomy tubes or other interventions to prevent complications.

\medterm{Ear, Nose, and Throat (ENT)}-- A medical specialty (otolaryngology) focused on the diagnosis and treatment of disorders affecting the ears, nose, throat, head, and neck. Conditions managed include hearing loss, sinusitis, tonsillitis, voice disorders, and head and neck cancers. Otolaryngologists also perform surgical procedures such as tonsillectomy, sinus surgery, and cochlear implantation.

\medterm{Earache (Otalgia)}-- Pain originating from the ear, a common symptom with diverse causes. Otalgia is classified as primary (originating within the ear itself) or secondary/referred (pain perceived in the ear but arising from another site). Primary causes: otitis media, otitis externa, Eustachian tube dysfunction, cholesteatoma, foreign body, cerumen impaction, barotrauma, mastoiditis. Referred causes: temporomandibular joint dysfunction, dental pathology (molar caries, abscess, impacted wisdom teeth), tonsillitis, pharyngitis, sinusitis, cervical spine disorders, trigeminal neuralgia, and malignancies of the pharynx or larynx. Evaluation: otoscopy, pneumatic otoscopy, tuning fork tests (Rinne, Weber), and assessment of surrounding structures. Management depends on etiology.

\medterm{Eardrum}-- A thin membrane separating the ear canal and middle ear.

\medterm{Earwax}-- Cerumen, a waxy secretion produced by the ceruminous glands in the outer third of the external auditory canal. Earwax lubricates and protects the skin of the ear canal, trapping dust, debris, and microorganisms to prevent them from reaching the tympanic membrane. Normal self-cleaning mechanism migrates wax laterally through epithelial migration. Impaction occurs when cerumen accumulates excessively, causing conductive hearing loss, itching, tinnitus, a feeling of fullness, or discomfort. Excessive use of cotton buds can cause impaction. Treatment: cerumenolytic drops (sodium bicarbonate, hydrogen peroxide, olive oil) or microsuction/ear irrigation. Irrigation is contraindicated with known tympanic membrane perforation, previous ear surgery, or active otitis media.

\medterm{Epistaxis}-- Epistaxis, commonly known as a nosebleed, is hemorrhage from the nasal mucosa that can range from minor, self-limited bleeding to severe, lifethreatening hemorrhage requiring emergency intervention, and it is a common presenting complaint in both primary care and emergency departments.

\medterm{Epley Maneuver}-- The Epley maneuver, also known as the canalith repositioning maneuver, is a therapeutic procedure used to treat posterior canal benign paroxysmal positional vertigo by repositioning displaced calcium carbonate crystals (otoconia) from the posterior semicircular canal back into the utricle through a series of sequential head and body position changes guided by the physician. The procedure begins with the patient seated, the head rotated forty-five degrees toward the affected ear, and the patient is rapidly laid back with the head hanging (the DixHallpike position), maintaining the head rotation for thirty to sixty seconds until the nystagmus resolves. The head is then rotated ninety degrees to the opposite side while maintaining the neck extension, held for thirty to sixty seconds, and the patient is then rolled onto their side with the nose pointing toward the floor, held for thirty to sixty seconds, before sitting upright. The maneuver relies on gravity to guide the otoconia through the semicircular canal and back into the utricle, with a success rate of approximately eighty percent after one to two treatments, and patients are advised to avoid lying flat on the affected side for several days afterward.

\medterm{Esophageal Foreign Body}-- Ingestion of a foreign object lodged in the esophagus. Most common in children (coins, batteries, toys) and adults with dentures (food bolus), psychiatric illness, or prisoners. Common sites: cricopharyngeus (C6, 60%), aortic arch (T4), and gastroesophageal junction. Symptoms: dysphagia, odynophagia, drooling, and foreign body sensation. Sharp objects: risk of perforation. Button batteries: require emergency removal due to rapid tissue necrosis and perforation. Diagnosis: X-ray neck/chest/abdomen (AP and lateral). Treatment: endoscopic removal (urgent for sharp, batteries, large objects, or complete obstruction). Glucagon IV for distal esophageal food impaction may facilitate passage. Bougienage for smooth coins in children.

\medterm{Ethmoidal Sinus}-- The ethmoidal sinuses, also known as ethmoidal air cells, are a collection of multiple small, thin-walled, air-containing cavities within the ethmoid bone between the orbits and above the nasal cavity, and they are divided into anterior and posterior groups based on their drainage sites. The anterior ethmoidal air cells, typically numbering six to twelve, drain through the ethmoidal infundibulum into the middle meatus, while the posterior ethmoidal air cells, typically numbering three to eight, drain into the superior meatus. The ethmoidal sinuses are separated from the orbits by the paper-thin lamina papyracea, and from the anterior cranial fossa by the fovea ethmoidalis (the roof of the ethmoid labyrinth), making complications of ethmoidal sinusitis potentially dan-

\medterm{External otitis}-- An infection or irritation of the outer ear, ear canal, or both. Also called swimmer’s ear.

\medterm{FESS}-- Functional endoscopic sinus surgery is a minimally invasive surgical technique that uses nasal endoscopes and specialized instruments to restore normal sinus drainage and ventilation by removing obstructive tissue and opening blocked sinus ostia, while preserving as much normal sinus mucosa and anatomy as possible. The procedure is indicated for chronic rhinosinusitis with or without nasal polyps that has failed medical therapy, recurrent acute sinusitis, and complications of sinusitis including mucocele and orbital complications. FESS involves the uncinectomy (removal of the uncinate process to access the maxillary sinus ostium), ethmoidectomy (removal of ethmoidal air cells), maxillary antrostomy (enlargement of the maxillary sinus ostium), frontal sinusotomy (opening of the frontal sinus drainage pathway), and sphenoidotomy (opening of the sphenoid sinus), depending on the extent of disease. The procedure is performed under general anesthesia using rigid endoscopes with angles of 0, 30, 45, and 70 degrees, and it offers the advantages of direct visualization, targeted tissue removal, preservation of normal anatomy, and faster recovery compared to traditional open sinus surgery.

\medterm{Frontal Sinus}-- The frontal sinuses are paired, air-filled cavities located within the frontal bone above the orbits and the root of the nose, separated by a bony septum that may be deviated to one side, and they are the last of the paranasal sinuses to develop, typically appearing around age seven and reaching full size by late adolescence. The frontal sinuses are irregular in

\medterm{Glossitis}-- Inflammation of the tongue. Causes: (1) Infection: bacterial (Streptococcus, Actinomyces), viral (HSV, Coxsackie--herpangina, HIV), fungal (Candida--oral thrush). (2) Nutritional deficiencies: iron (atrophic glossitis, smooth red tongue), B12 (magenta tongue), riboflavin (purple-red tongue with fissures), niacin (pellagra, beefy red tongue), folate. (3) Autoimmune: geographic tongue (benign migratory glossitis), lichen planus, pemphigus, Behçet disease. (4) Trauma: burns (hot food/drink), bite injury, ill-fitting dentures or orthodontic appliances. (5) Allergic: contact stomatitis (toothpaste, mouthwash, food). (6) Systemic: amyloidosis, sarcoidosis, Crohn disease. Symptoms: pain, burning, swelling, redness, altered taste, difficulty eating. Diagnosis: clinical examination, nutritional bloods, swab for culture, biopsy if persistent. Treatment: directed at cause (antifungals for Candida, B12 injections for pernicious anaemia, avoidance of irritants).

\medterm{Glottis}-- The opening between the vocal folds (true vocal cords) in the larynx, through which air passes during respiration and phonation. Boundaries: anterior (anterior commissure, attachment of vocal folds to thyroid cartilage), posterior (arytenoid cartilages and posterior commissure). The glottis has three functional states: (1) Open (respiration): vocal folds abducted by posterior cricoarytenoid muscles, creating a triangular opening. (2) Closed (swallowing, valsalva): vocal folds adducted, protecting the lower airway. (3) Oscillating (phonation): vocal folds vibrate as air passes through, producing sound. Glottic aperture size directly affects airway resistance and voice production. Glottic oedema (angioedema, anaphylaxis) causes stridor and airway obstruction.

\medterm{Hair Cells}-- The sensory receptor cells of the inner ear that convert mechanical sound vibrations into electrical neural signals, and they are named for the bundle of stereocilia projecting from their apical surface. The inner hair cells, arranged in a single row along the length of the cochlea, are the primary auditory receptors, with approximately 95 percent of the auditory nerve fibers innervating them, and they are responsible for transmitting sound information to the brain. The outer hair cells, arranged in three rows, function as cochlear amplifiers through a unique property called electromotility, in which the cells change length in response to electrical stimulation, actively amplifying the motion of the basilar membrane and enhancing the ear’s sensitivity to quiet sounds and its ability to distinguish between different frequencies. Hair cells are highly susceptible to damage from loud noise exposure, ototoxic medications (including aminoglycoside antibiotics, cisplatin, and loop diuretics), aging, and genetic mutations, and once damaged in mammals, they do not regenerate, leading to permanent sensorineural hearing loss.

\medterm{Hay Fever (Allergic Rhinitis)}-- An IgE-mediated, type I hypersensitivity reaction to airborne allergens (pollen, dust mites, mould spores, animal dander). Hay fever: seasonal allergic rhinitis (caused by grass, tree, and weed pollens), while perennial allergic rhinitis occurs year-round (house dust mites, pet dander, mould). Pathophysiology: allergen cross-links IgE on sensitised mast cells in nasal mucosa, triggering degranulation (histamine, leukotrienes, prostaglandins) within minutes (early phase: sneezing, itching, rhinorrhoea, nasal congestion). Late phase (4-6 hours): eosinophil and basophil recruitment, causing persistent congestion and hyperreactivity. Symptoms: sneezing paroxysms, anterior rhinorrhoea (clear watery discharge), nasal congestion, post-nasal drip, nasal pruritus, palatal/ear itching, conjunctival symptoms (itchy, red, watery eyes, periorbital oedema). Associated conditions: asthma (80% of asthmatics have rhinitis, unified airway concept), eczema, sinusitis, nasal polyps, otitis media with effusion. Diagnosis: clinical (symptom pattern, seasonal/perennial triggers), allergen-specific IgE testing (skin prick test or serum specific IgE, RAST). Management: (1) Allergen avoidance (encasing mattresses and pillows, HEPA filters, avoiding outdoor activity during high pollen counts, saline nasal irrigation). (2) Pharmacological: oral/nasal antihistamines (cetirizine, loratadine, fexofenadine), intranasal corticosteroids (beclomethasone, fluticasone, mometasone--first-line for moderate-severe intermittent or persistent disease), leukotriene receptor antagonists (montelukast--adjunctive with asthma), intranasal mast cell stabilisers (cromoglycate), oral decongestants (pseudoephedrine--short-term only). (3) Immunotherapy: allergen-specific subcutaneous (SCIT) or sublingual (SLIT) immunotherapy, modifies the immune response (induces IgG-blocking antibodies, regulatory T cells, shifts from Th2 to Th1), highly effective for pollen and dust mite allergies.

\medterm{Hearing aid}-- An electronic device worn in or behind the ear by people with hearing problems that makes sounds louder.

\medterm{Hearing Aids and Devices}-- Electronic devices worn in or behind the ear that amplify sound for individuals with hearing loss. Modern hearing aids include digital signal processing, directional microphones, and connectivity to smartphones and other audio sources. Cochlear implants are surgically implanted devices that directly stimulate the auditory nerve for those with severe-to-profound hearing loss who do not benefit from conventional hearing aids.

\medterm{Hearing Loss}-- Reduced ability to hear, classified as conductive (outer/middle ear pathology: cerumen impaction, otitis media, otosclerosis, tympanic membrane perforation), sensorineural (inner ear or auditory nerve: presbycusis, noise exposure, ototoxic drugs, Meniere disease, acoustic neuroma), or mixed. Presents with difficulty understanding speech, especially in background noise. Diagnosis: audiometry (pure tone average, speech discrimination), tympanometry. Treatment: hearing aids, cochlear implants for severe-to-profound SNHL, surgical correction for conductive causes. Prevention: noise protection, ototoxic medication monitoring.

\medterm{Impedance hearing testing}-- A test that sends sound waves to the eardrum to determine if a problem in the middle ear is causing hearing loss.

\medterm{Incus}-- The incus, commonly known as the anvil, is the middle ossicle of the three bones in the middle ear, articulating with the malleus at the incudomalleolar joint and with the stapes at the incudostapedial joint. It is the largest of the three ossicles and is shaped like a premolar tooth, with a body and two processes: the long process, which descends parallel to the manubrium of the malleus and articulates with the stapes via the lenticular process, and the short process, which projects posteriorly and is attached to the posterior wall of the middle ear by the posterior incudal ligament. The incus functions as a lever in the ossicular chain, amplifying sound vibrations from the tympanic membrane to the oval window with an amplification factor of approximately 1.3, and combined with the area ratio between the tympanic membrane and oval window, provides a total sound pressure amplification of approximately 22-fold. The long process of the incus is particularly vulnerable to ischemic necrosis due to its tenuous blood supply, and erosion of this structure is a common complication of chronic otitis media and cholesteatoma.

\medterm{Inner ear}-- The deepest part of the ear, consisting of the cochlea and the labyrinth.

\medterm{Kiesselbach Plexus}-- Kiesselbach plexus, also known as Little area, is a highly vascularized region located on the anterior nasal septum where five arterial anastomoses converge, making it the most common site of epistaxis (nosebleed), accounting for approximately ninety

\medterm{Labyrinth}-- The inner ear. It contains the cochlea, which is responsible for hearing, as well as structures that are needed for balance.

\medterm{Labyrinthitis}-- Inflammation of the labyrinth, the inner ear structure containing the cochlea and vestibular apparatus, resulting in both auditory and vestibular symptoms including sudden, severe vertigo, sensorineural hearing loss, tinnitus, nausea, and vomiting. The condition is distinguished from vestibular neuritis, which involves inflammation of the vestibular nerve without cochlear involvement, by the presence of hearing loss. Labyrinthitis may be caused by viral infections (the most common cause, often following upper respiratory infections), bacterial infections (typically spreading from middle ear infections or meningitis), autoimmune diseases, or ototoxic medications. The viral form is generally self-limited, with gradual improvement over days to weeks, though some patients experience persistent vestibular symptoms and residual hearing loss. Bacterial labyrinthitis is more severe and can result in permanent hearing loss and vestibular dysfunction, requiring aggressive antibiotic treatment and potentially surgical drainage of the middle ear. Treatment includes vestibular suppressants for acute vertigo, antiemetics, and vestibular rehabilitation therapy for persistent symptoms.

\medterm{Laryngitis}-- Inflammation of the larynx and vocal cords that presents with hoarseness, voice loss, sore throat, dry cough, and a sensation of a lump in the throat (globus sensation), and it may be acute or chronic. Acute laryngitis is most commonly caused by viral upper respiratory infections, and it is usually self-limited, resolving within one to two weeks, though it may be accompanied by fever, malaise, and cervical lymphadenopathy. Chronic laryngitis lasting more than three weeks may be caused by voice abuse, laryngopharyngeal reflux, chronic sinusitis with postnasal drip, smoking, environmental irritants, allergic rhinitis, and certain medications, and it can lead to chronic hoarseness and permanent vocal cord changes if not addressed. The diagnosis is typically clinical, based on history and laryngoscopy findings of vocal cord erythema, edema, and sometimes hemorrhagic spots. Treatment of acute laryngitis includes voice rest, hydration, humidification, and avoidance of irritants, while chronic laryngitis requires identification and treatment of the underlying cause.

\medterm{Laryngoscopy}-- A medical procedure used to visualise the larynx (voice box) and adjacent structures of the throat. Types: (1) Indirect laryngoscopy: the simplest technique, using a laryngeal mirror angled at 120 degrees to reflect the laryngeal image, with a headlight or external light source. (2) Flexible nasolaryngoscopy (or transnasal flexible laryngoscopy): a thin, flexible fibreoptic endoscope passed through the nostril to visualise the nasal cavity, pharynx, and larynx. Performed under topical anaesthesia (lidocaine spray, decongestant). The patient is awake and phonates to assess vocal fold movement. Used for evaluating hoarseness, stridor, dysphagia, laryngeal masses, and vocal cord paralysis. (3) Direct laryngoscopy: performed under general anaesthesia using a rigid laryngoscope (e.g., Macintosh or Miller blade for intubation; suspension laryngoscope for surgery), providing a direct view of the larynx and subglottis. Used for biopsy, removal of laryngeal lesions (polyps, nodules, papillomas), laser surgery, and foreign body removal. (4) Video laryngoscopy: uses a camera at the tip of a rigid laryngoscope blade, providing an indirect view of the glottis on a screen. Improves glottic visualisation in difficult airways (e.g., Cormack-Lehane grade 3-4, obesity, limited neck mobility). Devices: GlideScope, McGrath, C-MAC, Airtraq.

\medterm{Larynx}-- The larynx, commonly known as the voice box, is a complex organ located in the anterior neck between the pharynx and the trachea that serves three primary functions: phonation (voice production), protection of the lower airway from aspiration during swallowing, and regulation of airflow during respiration. The laryngeal skeleton consists of several cartilages including the thyroid cartilage (the largest, commonly known as the Adam’s apple), the cricoid cartilage (the only complete cartilaginous ring in the airway), the paired arytenoid cartilages (which articulate with the vocal cords), and the epiglottis (which covers the laryngeal inlet during swallowing). The larynx contains two pairs of folds: the true vocal cords (vocal folds), which are responsible for phonation and are covered by stratified squamous epithelium, and the false vocal cords (vestibular folds), which are located superior to the true cords and play a role in airway protection. The larynx is innervated by the superior laryngeal nerve (sensory to the supraglottic larynx) and the recurrent laryngeal nerve (motor to all intrinsic laryngeal muscles except the cricothyroid), and injury to these nerves,

\medterm{Malleus}-- The malleus, commonly known as the hammer, is the largest of the three ossicles in the middle ear and is the most laterally positioned, with its handle (manubrium) embedded in the fibrous layer of the tympanic membrane and its head articulating with the incus at the incudomalleolar joint. The malleus is an irregularly shaped bone approximately 8 to 9 millimeters in length, consisting of a head, neck, anterior process, lateral process, and manubrium, and it transmits vibrations from the tympanic membrane to the incus and subsequently to the stapes and oval window. The tensor tympani muscle attaches to the handle of the malleus and contracts reflexively in response to loud sounds, increasing the stiffness of the ossicular chain and reducing the transmission of potentially damaging low-frequency sounds to the inner ear. During middle ear infections, the malleus may become eroded by cholesteatoma or chronic inflammation, leading to conductive hearing loss, and its assessment during otoscopy is important for evaluating middle ear pathology.

\medterm{Mastoidectomy}-- A surgical procedure that involves the removal of diseased mastoid air cells and bone from the mastoid portion of the temporal bone, performed for chronic otitis media with cholesteatoma, mastoiditis, or as a component of surgery for acoustic neuroma and other skull base tumors. The canal wall up (closed) mastoidectomy technique preserves the bony wall between the external auditory canal and the mastoid cavity, maintaining the normal anatomy and eliminating the need for ongoing mastoid cavity care, while the canal wall down (open) mastoidectomy removes the bony wall, creating a common cavity between the external auditory canal and the mastoid, which is indicated for extensive cholesteatoma that cannot be completely removed through a canal wall up approach. The canal wall down technique results in a larger mastoid cavity that requires periodic cleaning and is prone to accumulation of debris and recurrent infection, but it provides better access and visualization for complete removal of cholesteatoma. The procedure is performed under general anesthesia using a postauricular

\medterm{Mastoiditis}-- Bacterial infection of the mastoid air cells, typically a complication of acute otitis media. Presents with postauricular pain, swelling, erythema, and tenderness, fever, and otorrhea. Diagnosis: CT temporal bone shows mastoid air cell opacification with bony destruction. Complications: subperiosteal abscess, Bezold abscess (pus tracking into sternocleidomastoid), facial nerve palsy, sigmoid sinus thrombosis, meningitis, and intracranial abscess. Treatment: IV antibiotics (ceftriaxone, vancomycin), myringotomy with tube placement, and mastoidectomy for refractory cases.

\medterm{Meatuses}-- The nasal meatuses are the air passages located beneath each corresponding turbinate on the lateral wall of the nasal cavity, and they serve as the drainage pathways for the paranasal sinuses and the nasolacrimal duct. The inferior meatus, located beneath the inferior turbinate, receives drainage from the nasolacrimal duct, which carries tears from the eye to the nasal cavity. The middle meatus, located beneath the middle turbinate, receives drainage from the maxillary sinus through the hiatus semilunaris, the frontal sinus through the frontonasal duct, and the anterior ethmoidal air cells through the ethmoidal infundibulum. The superior meatus, located beneath the superior turbinate, receives drainage from the posterior ethmoidal air cells. The sphenoethmoidal recess, located above and behind the superior turbinate, receives drainage from the sphenoid sinus. Understanding the drainage pathways of the paranasal sinuses through the meatuses is essential for diagnosing and treating sinusitis, as obstruction of these narrow drainage pathways can

\medterm{Meniere Disease}-- An inner ear disorder characterized by recurrent episodes of vertigo lasting 20 minutes to 12 hours, with fluctuating sensorineural hearing loss, tinnitus, and aural fullness. Due to endolymphatic hydrops. Diagnosis: clinical criteria; audiometry shows low-frequency SNHL during attacks. Treatment: low-sodium diet, diuretics (hydrochlorothiazide), intratympanic steroids, intratympanic gentamicin for vertigo control, and surgical interventions (endolymphatic sac decompression, vestibular nerve section) for refractory cases.

\medterm{Middle Ear}-- The middle ear, also known as the tympanic cavity, is an air-filled, mucosa-lined space within the temporal bone that lies between the tympanic membrane laterally and the medial wall of the middle ear (abyrinthine wall) medially, and it contains the three ossicles (malleus, incus, and stapes) that transmit sound vibrations from the tympanic membrane to the oval window of the inner ear. The middle ear is divided into three compartments: the mesotympanum (directly behind the tympanic membrane), the epitympanum or attic (above the level of the tympanic membrane), and the hypotympanum (below the level of the tympanic membrane). The middle ear is connected to the nasopharynx by the Eustachian tube, which equalizes air pressure on both sides of the tympanic membrane, and it is lined by a mucous membrane that is continuous with the lining of the Eustachian tube and the pharynx. The

\medterm{Mixed Hearing Loss}-- A combination of both conductive and sensorineural hearing loss occurring in the same ear, resulting in a hearing deficit that involves impaired sound transmission through the outer or middle ear in addition to cochlear or neural dysfunction. The audiometric pattern shows both an air-bone gap indicating the conductive component and elevated bone conduction thresholds indicating the sensorineural component, and the total hearing loss is the sum of both components. Common causes include chronic otitis media or cholesteatoma with superimposed presbycusis, otosclerosis with cochlear involvement, and noise exposure with concurrent middle ear pathology. The management of mixed hearing loss requires addressing both components: treating the conductive component medically or surgically when possible (such as repairing a tympanic membrane perforation or reconstructing the ossicular chain) and managing the sensorineural component with hearing amplification or cochlear implantation as indicated.

\medterm{Nasal Cavity}-- The nasal cavity is the paired, air-filled space located between the nasal septum medially and the lateral nasal wall, extending from the nares (nostrils) anteriorly to the choanae (posterior openings into the nasopharynx) posteriorly, and it serves three primary functions: warming, humidifying, and filtering inspired air; olfaction (smell); and resonance of the voice. The lateral nasal wall contains three scroll-like bony projections called turbinates (superior, middle, and inferior), which increase the surface area of the nasal mucosa and create turbulent airflow to optimize air conditioning. Beneath each turbinates lies a corresponding meatus: the superior, middle, and inferior meatuses, which receive drainage from the paranasal sinuses. The nasal septum, composed of both cartilage and bone, divides the nasal cavity into two halves and is lined by a

\medterm{Nasal Polyps}-- Benign, pedunculated, edematous outgrowths of the sinus and nasal mucosa that arise most commonly from the ethmoidal sinuses and protrude into the nasal cavity through the middle meatus, causing nasal obstruction, anosmia (loss of smell), and rhinorrhea. The polyps are composed of edematous stroma with a sparse inflammatory infiltrate including eosinophils, lymphocytes, and plasma cells, and they are associated with type 2 inflammation, elevated tissue and serum IgE levels, and aspirin sensitivity in approximately thirty to forty percent of patients (Samter triad: aspirin sensitivity, asthma, and nasal polyposis). Nasal polyps are more common in adults, particularly males over the age of forty, and they are associated with chronic rhinosinusitis, asthma, aspirin sensitivity, cystic fibrosis, and primary ciliary dyskinesia.

\medterm{Nasal Septum}-- The nasal septum is the midline structure that divides the nasal cavity into two halves, consisting of an anterior cartilaginous portion (the septal cartilage or quadrilateral cartilage), a posterior bony portion (the perpendicular plate of the ethmoid bone and the vomer), and the membranous portion (the columella) at the anterior nasal opening. The septum provides structural support for the nose, divides the nasal airflow, and supports the nasal mucosa, which is richly vascularized by branches of the internal and external carotid arteries. Deviation of the nasal septum is extremely common, with virtually all individuals having some degree of asymmetry, and significant deviation can cause nasal obstruction, impaired sinus drainage, recurrent

\medterm{Nasopharyngeal Carcinoma}-- A malignancy arising from the epithelial lining of the nasopharynx, and it has a unique geographical and ethnic distribution, with highest incidence in southern China, Southeast Asia, and North Africa. The tumor is strongly associated with Epstein-Barr virus infection, which is present in virtually all cases of undifferentiated nasopharyngeal carcinoma, and dietary factors including salt-preserved fish and other nitrosaminecontaining foods. The tumor typically arises in the lateral wall of the nasopharynx (fossa of Rosenmuller) and often presents with a painless neck mass (cervical lymphadenopathy), nasal obstruction, epistaxis, hearing loss from Eustachian tube obstruction, and cranial nerve palsies. Diagnosis is confirmed by endoscopic examination and biopsy, and staging involves MRI of the nasopharynx and neck and PET-CT. Nasopharyngeal carcinoma is highly radiosensitive, and the primary treatment is radiation therapy combined with chemotherapy for locally advanced disease, with chemotherapy regimens typically including cisplatin and 5-fluorouracil.

\medterm{Neck Cancer}-- Human papillomavirus, particularly HPV-16, is an increasingly important etiological factor in head and neck squamous cell carcinoma, particularly oropharyngeal cancer, and it is responsible for the rising incidence of oropharyngeal cancer in younger patients without traditional risk factors of tobacco and alcohol use. HPV-positive oropharyngeal cancer has distinct epidemiological, clinical, and molecular characteristics compared to HPV-negative disease, including a younger age of presentation, better performance status, smaller primary tumors with larger cervical lymph node metastases, and significantly improved prognosis, with five-year survival rates of approximately eighty to ninety percent compared

\medterm{Non-Allergic Rhinitis}-- Non-allergic rhinitis, also known as vasomotor rhinitis, is a chronic nasal condition characterized by nasal symptoms including congestion, rhinorrhea, and sneezing that are not caused by an IgE-mediated allergic response, and it is distinguished from allergic rhinitis by the absence of allergen-specific IgE, negative allergy testing, and the absence of the systemic symptoms typically associated with allergic rhinitis. The pathogenesis involves dysfunction of the autonomic nervous system regulation of the nasal mucosal blood vessels and glands, leading to exaggerated vasodilation and increased mucus production in response to environmental triggers

\medterm{Oral Leukoplakia}-- A white patch or plaque on the oral mucosa that cannot be scraped off and cannot be characterized as any other diagnosable disease, considered a potentially malignant disorder. Risk factors: tobacco use (most common), alcohol, and betel nut chewing. Most common sites: buccal mucosa, alveolar ridge, and tongue. Dysplasia (mild, moderate, severe) found on biopsy in 10-40% of cases. Progression to invasive squamous cell carcinoma: ~1% per year for non-dysplastic, higher with dysplasia. Diagnosis: biopsy for histology. Treatment: excision or laser ablation for dysplastic lesions; observation for non-dysplastic. Risk factor modification essential. Long-term surveillance required.

\medterm{Ossicles}-- Three bones in the middle ear that move in response to sound vibrations.

\medterm{Otic capsule}-- The otic capsule is the dense bony shell (otic capsule bone) that surrounds and protects the inner ear structures, including the cochlea, vestibule, and semicircular canals, forming the hardest bone in the human body. It develops from the otic placode and undergoes endochondral ossification, and unlike other temporal bone components, it undergoes very little remodeling after development, making it resistant to otosclerosis but susceptible to otic capsule dehiscence (superior semicircular canal dehiscence syndrome), in which thinning or absence of the bone over the superior semicircular canal creates a third window phenomenon, causing sound- or pressure-induced vertigo (Tullio phenomenon), conductive hearing loss, and autophony.

\medterm{Otitis Externa}-- Otitis externa, commonly known as swimmer’s ear, is an infection or inflammation of the external

\medterm{Otitis media}-- A group of inflammatory conditions of the middle ear, most commonly caused by viral or bacterial infection, with Streptococcus pneumoniae, Haemophilus influenzae, and Moraxella catarrhalis being the most frequent bacterial pathogens. It is classified into acute otitis media (AOM), characterized by rapid-onset otalgia, fever, and bulging erythematous tympanic membrane with middle ear effusion; otitis media with effusion (OME), in which non-purulent fluid persists in the middle ear without signs of acute infection, often causing conductive hearing loss and speech delay in children; and chronic suppurative otitis media (CSOM), defined by persistent otorrhea through a tympanic membrane perforation for more than 6 weeks. Eustachian tube dysfunction is the primary predisposing factor, as impaired drainage and ventilation of the middle ear allow pathogens to ascend from the nasopharynx. Treatment of AOM includes observation or antibiotics (amoxicillin 80-90 mg/kg/day for 10 days in children) depending on age and severity, while OME may require tympanostomy tube placement if effusion persists beyond 3 months with hearing loss. Complications include tympanic membrane perforation, mastoiditis, cholesteatoma, and intracranial spread.

\medterm{Otolaryngology (ENT)}-- The surgical specialty concerned with the diagnosis and treatment of disorders of the ear, nose, throat, head, and neck. Otolaryngologists (ENT surgeons) manage conditions ranging from otitis media and sinusitis to head and neck cancer, thyroid disease, and voice disorders. Subspecialties: otology/neurotology (ear disorders, hearing loss, balance), rhinology (sinus disease, anterior skull base), laryngology (voice, swallowing), head and neck oncology (cancer surgery including neck dissection, laryngectomy, thyroidectomy), paediatric ENT (cleft palate, cochlear implants), and facial plastic and reconstructive surgery. Training in the UK requires 2 years of core surgical training followed by 6 years of specialty training (ST3--ST8).

\medterm{Otology}-- The branch of medicine and surgery concerned with the study and treatment of ear disorders. Otology encompasses hearing loss (conductive, sensorineural, mixed), tinnitus, vertigo and balance disorders (Meniere disease, benign paroxysmal positional vertigo, vestibular neuritis), otitis media, otosclerosis, cholesteatoma, and ear trauma. Diagnostic tools: pure-tone audiometry, tympanometry, otoacoustic emissions, auditory brainstem response (ABR), and CT/MRI of the temporal bone. Surgical procedures include myringotomy, tympanoplasty, mastoidectomy, stapedectomy, cochlear implantation, and ossiculoplasty. Neurotology is a subspecialty addressing the interface between the ear and brain (acoustic neuroma, CSF leak repair, skull base surgery).

\medterm{Otorrhoea}-- Discharge from the ear. Classification: acute otorrhoea (<6 weeks) and chronic otorrhoea (>6 weeks). Causes: acute otitis media with tympanic membrane perforation (purulent discharge, often in children), otitis externa (watery or mucoid discharge, ear canal inflammation), chronic suppurative otitis media (persistent purulent discharge through a non-healing perforation, associated with cholesteatoma), cerebrospinal fluid otorrhoea (clear, watery, \(\beta\)-2 transferrin positive, indicates skull base fracture or surgical complication), and foreign body. Otorrhoea with pain, fever, and mastoid tenderness raises concern for acute mastoiditis. Management: topical antibiotic drops (ciprofloxacin—first-line, avoids ototoxicity of aminoglycosides), aural toilet (microsuction), treatment of underlying cause; for CSF otorrhoea, neurosurgical repair is required.

\medterm{Otosclerosis}-- An abnormal bony growth in the middle ear that causes hearing loss.

\medterm{Otoscopy}-- A clinical examination technique used to visualize the external auditory canal and tympanic membrane, and it is an essential component of the evaluation of ear complaints including hearing loss, ear pain, otorrhea, and tinnitus. Pneumatic otoscopy involves the use of a specialized otoscope with a sealed speculum that allows insufflation of air into the ear canal while observing the tympanic membrane, assessing the mobility of the membrane in response to pressure changes, which is reduced in otitis media with effusion and absent in middle ear fluid. The normal tympanic membrane appears as a translucent, pearly gray, cone-shaped membrane with a visible light reflex in the anteroinferior quadrant, and the handle of the malleus and the umbo (the central point of the tympanic membrane) are visible landmarks. Abnormal findings include bulging, erythema, effusion, perforation, retraction, cholesteatoma, and tympanostomy tubes, and proper otoscopic technique requires pulling the pinna superiorly and posteri-

\medterm{Outer ear}-- The external part of the ear, as well as the external ear canal and the eardrum.

\medterm{Ozaena}-- A chronic condition characterised by atrophic rhinitis with crusting, foul-smelling nasal discharge (fetor), and a sensation of nasal obstruction despite a patent airway. Ozaena is caused by progressive atrophy of the nasal mucosa and underlying bone, often associated with Klebsiella ozaenae infection. It is more common in women, typically presents in middle age, and is more prevalent in tropical and subtropical regions. The nasal mucosa is pale, dry, and crusted; removal of crusts reveals a widely patent nasal cavity. Atrophy may extend to the olfactory epithelium, causing anosmia. Treatment: nasal saline irrigation, topical antibiotics (neomycin, gentamicin), oral antibiotics (ciprofloxacin, rifampicin for Klebsiella), oestrogen nasal spray (improves mucosal blood flow), and surgical closure of the nasal cavity (Young procedure) in severe refractory cases.

\medterm{Paranasal Sinuses}-- The paranasal sinuses are four pairs of air-filled, mucosa-lined cavities within the facial bones that communicate with the nasal cavity through narrow drainage pathways called ostia, and they serve multiple functions including lightening the weight of the skull, humidifying and warming inspired air, producing mucus for nasal clearance, contributing to vocal resonance, and providing structural support for the face. The maxillary sinuses, the largest of the paranasal sinuses, are located in the maxillary bones beneath the orbits and lateral to the nasal cavity, with their ostia opening into the middle meatus. The frontal sinuses, located in the frontal bone above the orbits, drain through the frontonasal duct into the middle meatus. The ethmoidal sinuses, consisting of multiple small air cells within the ethmoid bone between the orbits, are divided into anterior and posterior groups, with the anterior cells draining into the middle meatus and the posterior cells draining into the superior meatus. The sphenoid sinuses, located in the sphenoid bone posterior to the ethmoidal sinuses, drain into the sphenoethmoidal recess. Sinusitis occurs when the ostia become obstructed, preventing mucus drainage and creating a favorable environment for bacterial infection.

\medterm{Parotitis}-- Inflammation of the parotid gland. Acute parotitis: most commonly caused by mumps virus (paramyxovirus—epidemic parotitis, typically in unvaccinated children, bilateral in 70%, resolves spontaneously over 7--10 days). Suppurative (bacterial) parotitis: Staphylococcus aureus (most common), streptococci, anaerobes—occurs in dehydrated, post-surgical, or immunocompromised patients; presents with painful, erythematous swelling, purulent discharge from Stensen duct, fever, and leukocytosis; treated with IV antibiotics (flucloxacillin), hydration, sialagogues, and surgical drainage if abscess forms. Chronic parotitis: recurrent episodes due to duct obstruction (sialolithiasis), Sjögren syndrome, sarcoidosis, or IgG4-related disease. Diagnosis: clinical, ultrasound, sialography, MRI. Complications of bacterial parotitis: parotid abscess, facial nerve palsy, extension to the parapharyngeal space. Mumps parotitis is prevented by MMR vaccination (two doses).

\medterm{Peritonsillar Abscess}-- Peritonsillar abscess, also known as quinsy, is a collection of pus in the peritonsillar space between the tonsillar capsule and the superior pharyngeal constrictor muscle, typically developing as a complication of acute tonsillitis, and it is the most common deep neck space infection. The condition presents with severe, unilateral sore throat, odynophagia, dysphagia, trismus (difficulty opening the mouth due to spasm of the pterygoid muscles), hot potato voice (muffled speech), and deviation of the uvula toward the contralateral side, with erythematous, bulging of the soft palate and tonsillar pillar on the affected side. The most common causative organisms

\medterm{Pharyngitis}-- Inflammation of the pharynx, the most common presentation of which is sore throat (odynophagia) with or without fever, and it is classified as bacterial or viral based on the causative organism, with viral pharyngitis accounting for approximately eighty to ninety percent of cases and bacterial pharyngitis (primarily Group A Streptococcus) accounting for approximately ten to twenty percent. Viral pharyngitis typically presents with a gradual onset of sore throat, cough, rhinorrhea, conjunctivitis, and hoarseness, and is caused by adenoviruses, rhinoviruses, coronaviruses, influenza, and

\medterm{Pleomorphic Adenoma}-- Pleomorphic adenoma, also known as benign mixed tumor, is the most common benign salivary gland tumor, accounting for approximately sixty to seventy percent of all parotid gland neoplasms and a significant proportion of minor salivary gland tumors, and it is characterized by a diverse histological appearance with a mixture of epithelial, myoepithelial, and mesenchymal-like components in a variable stroma that may be mucoid, chondroid, or osseous. The tumor typically presents as a slow-growing, painless, firm, mobile mass in the parotid gland, most commonly in the superficial lobe, and it is more common in women (female-to-male ratio of approximately 2:1), with a peak incidence in the fourth to

\medterm{Pneumatic Otoscopy}-- A diagnostic technique that combines visualization of the tympanic membrane with assessment of its mobility in response to positive and negative air pressure changes, providing information about the presence of middle ear fluid and the condition of the Eustachian tube. The procedure uses a pneumatic otoscope with a soft rubber bulb that generates positive and negative pressure, and the examiner observes the tympanic membrane for movement in response to these pressure changes. Normal tympanic membranes move freely with pressure changes, while middle ear effusion reduces or eliminates mobility, and the degree of reduced mobility correlates with the presence and volume of middle ear fluid. Pneumatic otoscopy is recommended as the primary diagnostic tool for acute otitis media in children, as it provides both structural and functional information about the middle ear, and it can detect middle ear effusion with a sensitivity of approximately ninety percent when performed by experienced examiners. The technique is inexpensive, non-invasive, and does not require sedation, making it ideal for use in young children.

\medterm{Presbycusis}-- Age-related hearing loss caused by the death of hair cells in the inner ear.

\medterm{Prosthesis}-- An artificial device such as a hearing aid, artificial joint, or dentures that substitutes for a missing body part.

\medterm{Pure Tone Audiometry}-- The standard behavioral test used to measure hearing sensitivity across the frequency range important for speech perception, typically from 250 to 8000 Hertz, by determining the lowest intensity level at which a pure tone can be detected approximately fifty percent of the time (the hearing threshold). Air conduction thresholds are measured using headphones or insert earphones, and bone conduction thresholds are measured using a bone vibrator placed on the mastoid process or forehead, and the comparison between air and bone conduction thresholds distinguishes between conductive, sensorineural, and mixed hearing loss. The audiogram plots the hearing thresholds in decibels hearing level (dB HL) against frequency in Hertz (Hz), with normal hearing defined as thresholds of 25 dB HL or better. The degree of hearing loss is classified as normal (0-25 dB HL), mild (26-40 dB HL), moderate (41-55 dB HL), moderately severe (56-70 dB HL), severe (71-90 dB HL), and profound (greater than 90 dB HL). Pure tone audiometry is the foundation of audiometric assessment and is essential for diagnosing hearing loss, determining its type and degree, and guiding treatment decisions.

\medterm{Quinsy}-- See Peritonsillar Abscess.

\medterm{Retropharyngeal Abscess}-- A collection of pus in the retropharyngeal space, located between the posterior pharyngeal wall and the prevertebral fascia. Most common in children aged 2-4 years, arising from suppurative lymphadenitis of the retropharyngeal lymph nodes that drain the nasopharynx, adenoids, and sinuses. In adults, it typically results from direct extension of pharyngeal infection (tonsillitis, dental infection) or from penetrating trauma (foreign body, instrumentation). Presents with fever, sore throat, neck pain, dysphagia, drooling, cervical lymphadenopathy, and neck stiffness. In children, the classic presentation includes refusal to feed, stridor, or respiratory distress. Diagnosis is confirmed by contrast-enhanced CT of the neck showing a hypodense collection with rim enhancement in the retropharyngeal space. Lateral neck X-ray may show prevertebral soft tissue swelling (greater than 7 mm at C2 or greater than 14 mm at C6). Treatment: airway management, IV antibiotics (e.g., clindamycin or ampicillin-sulbactam), and surgical drainage via transoral incision or transcervical approach for larger abscesses or those not responding to antibiotics. Complications include airway obstruction, mediastinitis, jugular vein thrombosis, and sepsis.

\medterm{Rhinitis}-- Inflammation of the nasal mucosa. Allergic rhinitis: IgE-mediated response to environmental allergens (pollen, dust mites, pet dander), presenting with sneezing, rhinorrhea, nasal congestion, and pruritus. Non-allergic rhinitis: vasomotor rhinitis (triggers: irritants, temperature changes, strong odors), infectious (viral URI), medication-induced (rhinitis medicamentosa from topical decongestant overuse), hormonal (pregnancy, hypothyroidism), and gustatory rhinitis (food-induced). Diagnosis: clinical, allergy testing. Treatment: intranasal corticosteroids (first-line), oral antihistamines, decongestants, saline irrigation, allergen immunotherapy for allergic rhinitis.

\medterm{Rhinoplasty}-- A surgical procedure to alter the shape, size, or function of the nose, performed for cosmetic reasons (aesthetic rhinoplasty—improving nasal appearance) or functional reasons (functional rhinoplasty—improving nasal airflow by correcting internal/external nasal valve collapse, deviated septum, or turbinate hypertrophy). Approaches: open rhinoplasty (transcolumellar incision—better exposure, preferred for complex cases) and closed (endonasal) rhinoplasty (incisions within the nostril—no external scar, less swelling). Components: hump reduction (dorsal bony and cartilaginous hump), osteotomies (narrowing the nasal bones), tip rhinoplasty (refinement of the nasal tip—cephalic trim, tip sutures, columellar strut graft), alar base reduction, and septoplasty (correction of septal deviation). Grafts: autologous cartilage (septal, auricular, costal) for structural support or contour augmentation. Post-operative: swelling and bruising (peaks at 48 hours, resolves over 1--2 weeks), nasal splint for 1 week, avoidance of glasses and strenuous activity for 4--6 weeks. Complications: bleeding, infection, nasal obstruction, asymmetry, saddle nose deformity (over-resection of dorsum), and need for revision (5--15%).

\medterm{Rinne Test}-- The Rinne test is a screening test that compares air conduction and bone conduction hearing using a vibrating tuning fork (typically 512 Hz) to help differentiate between conductive and sensorineural hearing loss. The test is performed by first placing the vibrating tuning fork on the mastoid process to assess bone conduction, asking the patient to indicate when the sound is no longer heard, and then immediately holding the tuning fork adjacent to the external auditory canal to assess air conduction, asking the patient to indicate whether the sound is still audible. A positive Rinne test (normal result) occurs when air conduction is heard longer than bone conduction, indicating normal middle ear function or sensorineural hearing loss. A negative

\medterm{Salivary Gland Tumors}-- Most salivary tumors are benign (80% in parotid, 50% in submandibular, 50% in minor glands). Benign: pleomorphic adenoma (most common overall), Warthin tumor (bilateral in 10%, smokers). Malignant: mucoepidermoid carcinoma (most common malignant), adenoid cystic carcinoma (perineural invasion), acinic cell carcinoma. Presents with painless, slow-growing mass. Malignant features: rapid growth, pain, facial nerve paralysis, skin fixation, and cervical lymphadenopathy. Diagnosis: fine needle aspiration (FNA), CT/MRI for staging. Treatment: superficial parotidectomy for benign; total parotidectomy with neck dissection for malignant, with postoperative radiation as needed.

\medterm{Semicircular Canals}-- The semicircular canals are three bony, fluid-filled, tube-like structures within the vestibular portion of the inner ear that detect rotational (angular) acceleration

\medterm{Septal Hematoma}-- Accumulation of blood between the perichondrium and the septal cartilage of the nose, typically following nasal trauma (nasal fracture). Presents with nasal obstruction, pain, and a bluish, boggy swelling of the nasal septum. An untreated septal hematoma causes avascular necrosis of the nasal cartilage (cartilage receives nutrition from perichondrium), leading to saddle nose deformity (septal perforation and collapse) within days. Diagnosis: nasal speculum exam. Treatment: immediate incision and drainage (small incision at the inferior aspect), packing to prevent re-accumulation, and antibiotics to prevent infection (septal abscess). Follow-up for re-accumulation.

\medterm{Sialolithiasis}-- Sialolithiasis, commonly known as salivary stones, is the formation of calcified concretions within the salivary gland ducts or parenchyma, most commonly affecting the submandibular gland (approximately eighty percent of cases) due to the long, dependent course of the Wharton duct, the alkaline pH of submandibular saliva, and its higher mucin and calcium content, followed by the parotid gland (approximately ten percent) and the sublingual gland (approximately five percent). The stones are composed primarily of calcium phosphate and calcium carbonate, with an organic matrix of glycoproteins and bacteria, and they range in size from a few millimeters to several centimeters. Patients typically present with intermittent, recurrent episodes of glandular swelling and pain that worsen during meals (when salivary flow is stimulated), and the stone may be palpable in the floor of the mouth (submandibular stones) or intraorally. Diagnosis is confirmed by plain radiography (submandibular stones are radiopaque in approximately eighty percent), ultrasound, or CT imaging. Treatment depends on stone size and location, with small stones managed with sialagogues and hydration, and larger or impacted stones managed with sialendoscopy, lithotripsy, or surgical removal.

\medterm{Sialorrhoea}-- See Drooling / Excessive Salivation.

\medterm{Sinusitis}-- Inflammation of the paranasal sinus mucosa. Acute sinusitis: viral (most common, resolves in <10 days) or bacterial (symptoms >10 days, double worsening, or severe: Streptococcus pneumoniae, Haemophilus influenzae, Moraxella catarrhalis). Chronic sinusitis: symptoms >12 weeks with radiographic evidence. Presents with nasal congestion, purulent rhinorrhea, facial pain/pressure, headache, and hyposmia. Diagnosis: clinical, CT sinus for chronic/recurrent. Treatment: analgesics, saline irrigation, intranasal corticosteroids; antibiotics for acute bacterial. Endoscopic sinus surgery for refractory chronic sinusitis.

\medterm{Sinusitis and Sinus Problems}-- Sinus problems refer to conditions affecting the paranasal sinuses, which are air-filled cavities in the skull that can become inflamed or infected. Acute sinusitis typically follows a viral upper respiratory infection and causes facial pain, nasal congestion, and headache, while chronic sinusitis persists for weeks or longer and may involve nasal polyps or allergic components. Treatment ranges from saline rinses and decongestants to intranasal corticosteroids, antibiotics for bacterial infection, and in refractory cases, endoscopic sinus surgery.

\medterm{Speech Discrimination}-- Speech discrimination testing, also known as word recognition scoring, measures the ability to understand speech at a comfortable listening level, and it is an important component of the audiometric evaluation that provides information about cochlear and neural function beyond what pure tone audiometry alone can reveal. The test typically uses a standardized list of phonetically balanced monosyllabic words presented at a suprathreshold level (usually 40 dB above the speech recognition threshold), and the patient repeats the words, with the percentage of words correctly identified representing the speech discrimination score. Normal speech discrimination scores range from ninety to one hundred percent, while reduced scores indicate cochlear pathology (particularly in patients with high-tone sensorineural hearing loss) or retrocochlear pathology (as in acoustic neuroma, where scores may be disproportionately

\medterm{Sphenoidal Sinus}-- The sphenoid sinuses are paired, air-filled cavities located within the body of the sphenoid bone, the most posterior of the paranasal sinuses, situated deep within the skull at the center of the cranial base, and they are in close proximity to critical neurovascular structures including the optic nerves, the internal carotid arteries, the cavernous sinuses, the pituitary gland, and the brainstem. The sphenoid sinuses vary considerably in size and pneumatization among individuals, and they drain through the sphenoethmoidal recess located above and posterior to the superior turbinate into the superior meatus.

\medterm{Stapedectomy}-- A microsurgical procedure performed to treat otosclerosis, a condition characterized by abnormal bone remodeling in the otic capsule that fixes the stapes footplate in the oval window, preventing its normal vibratory movement and causing progressive conductive hearing loss. The procedure involves creating a small opening in the tympanic membrane, elevating the tympanomeatal flap to visualize the middle ear, identifying the fixed stapes, and either removing the entire stapes footplate (stapedectomy) or creating a small fenestration in the footplate (stapedotomy) and inserting a prosthetic device (typically a titanium or fluoroplastic piston) that connects the incus to the oval window, bypassing the fixed stapes and restoring the ossicular chain’s ability to transmit sound vibrations. The procedure is performed under local or general anesthesia, and the success rate for hearing improvement is approximately ninety to ninety-five percent, with closure of the air-bone gap to within 10 decibels in most patients. Complications include sensorineural hearing loss (from perilymph leak or inner ear trauma), vertigo, taste disturbance, and tympanic membrane perforation.

\medterm{Stapes}-- The stapes, commonly known as the stirrup, is the smallest bone in the human body and the most medially positioned ossicle in the middle ear, articulating with the incus at the incudostapedial joint and transmitting sound vibrations to the inner ear through its footplate, which is attached to the oval window by the annular ligament. The stapes consists of a head, two crura (anterior and posterior), and a footplate approximately 3.2 millimeters by 1.4 millimeters, and it acts as a piston, converting the lever action of the malleus and incus into a piston-like movement that transmits sound energy into the fluid-filled inner ear. The stapedius

\medterm{Stomatitis}-- Inflammation of the oral mucosa, presenting as erythema, ulceration, pain, and difficulty eating. Causes: infections (viral—herpes simplex (herpetic gingivostomatitis), Coxsackie (hand-foot-and-mouth disease), herpes zoster; bacterial—acute necrotising ulcerative gingivitis (Vincent angina); fungal—Candida albicans (oral thrush, common in immunocompromised, infants, denture wearers, after antibiotics/inhaled corticosteroids)), physical trauma (biting, burns, ill-fitting dentures), chemical irritation (aspirin burn, alcohol, smoking), aphthous ulcers (recurrent aphthous stomatitis—minor, major, herpetiform; associated with stress, trauma, nutritional deficiencies, Crohn disease, Behçet disease), autoimmune (pemphigus vulgaris, mucous membrane pemphigoid, lichen planus), drug-induced (methotrexate, NSAIDs, chemotherapy—mucositis), and nutritional deficiencies (B vitamins, iron, folate, zinc). Diagnosis: clinical, swab for culture/PCR, patch testing, and biopsy. Treatment: treat underlying cause; symptomatic—topical analgesics (lidocaine), antiseptic mouthwashes (chlorhexidine), corticosteroids (topical or systemic), antifungals for Candida, antivirals for herpes.

\medterm{Sublingual Gland}-- The sublingual glands are the smallest of the major salivary glands, located in the floor of the mouth beneath the tongue, and they produce primarily mucous saliva that is delivered into the oral cavity through multiple small ducts (ducts of Rivinus) that open along the sublingual fold, and through the major sublingual duct (Bartholin duct), which joins the Wharton duct of the submandibular gland. The sublingual glands are rarely affected by sialolithiasis due to the short, wide ducts, and they are more commonly involved in autoimmune conditions including

\medterm{Thyroid Cartilage}-- The thyroid cartilage is the largest of the laryngeal cartilages, consisting of two laminae that fuse anteriorly at an angle of approximately ninety degrees in males and one hundred twenty degrees in females, creating the laryngeal prominence (commonly known as the Adam’s apple), which is more prominent in males due to the acute angle and the effects of testosterone during puberty. The thyroid cartilage forms the anterior and lateral walls of the larynx and provides structural support and protection for the vocal cords and the intrinsic laryngeal muscles. The superior and inferior horns of the thyroid cartilage articulate with the hyoid bone and the cricoid cartilage, respectively, and the cricothyroid muscle, the only intrinsic laryngeal muscle innervated by the external branch of the superior laryngeal nerve, attaches to the thyroid cartilage and acts to tense the vocal cords, raising the pitch of the voice. The thyroid cartilage is prone to fracture in blunt neck trauma, and thyroid cartilage fractures can cause airway compromise, subcutaneous emphysema, and hoarseness.

\medterm{Thyroid Nodule}-- A thyroid nodule is a discrete lesion within the thyroid gland that is palpably or radiographically distinct from the surrounding thyroid parenchyma, and it is a common clinical finding, present in approximately five percent of palpable neck examinations and up to fifty percent of individuals on ultrasound imaging, with the vast majority being benign. The evaluation of thyroid nodules focuses on identifying the small minority that represent thyroid carcinoma, and it involves ultrasound assessment of nodule characteristics (size, composition, echogenicity, margins, shape, and calcifications), fine-needle aspiration biopsy of suspicious nodules, and thyroid function testing. The Thyroid Imaging Reporting

\medterm{TI-RADS}-- The Thyroid Imaging Reporting and Data System is a standardized ultrasound-based risk stratification system developed by the American College of Radiology to evaluate thyroid nodules and guide recommendations for fine-needle aspiration biopsy, reducing unnecessary biopsies of benignappearing nodules while ensuring adequate sampling of suspicious nodules. The TI-RADS scoring system assigns points based on five ultrasound characteristics: composition (solid, mixed cystic-solid, predominantly cystic, or spongiform), echogenicity (anechoic, hypoechoic, isoechoic, or markedly hypoechoic), shape (wider-than-tall or taller-thanwide), margin (smooth, ill-defined, lobulated, or extra-thyroidal extension), and echogenic foci (none, macrocalcifications, rim calcifications, or microcalcifications). The total TI-RADS score determines the category from TR1 (benign, no biopsy needed) to TR5 (highly suspicious, biopsy recommended for nodules one centimeter or larger), and the system provides evidence-based size thresholds for biopsy recommendations at each TI-RADS level, optimizing the balance between cancer detection and unnecessary procedures.

\medterm{Tinnitus}-- The perception of sound in the absence of an external acoustic stimulus, commonly described as ringing, buzzing, hissing, or roaring in the ears, and it affects approximately fifteen to twenty percent of the global population, with a significant impact on quality of life in severe cases. Subjective tinnitus, the most common form, is perceived only by the patient and is most commonly associated with sensorineural hearing loss, noise ex-

\medterm{Tonsillectomy}-- The surgical removal of the palatine tonsils, and it is one of the most commonly performed surgical procedures in children, traditionally indicated for recurrent acute tonsillitis (seven episodes in one year, five per year for two years, or three per year for three years), chronic tonsillitis, and obstructive sleep apnea caused by tonsillar hypertrophy. The procedure is performed under general anesthesia, and the most commonly used techniques include cold dissection with electrocautery, bipolar diathermy, intracapsular tonsillectomy (partial tonsillectomy using a microdebrider or coblation), and total tonsillectomy. The most common complications include postoperative hemorrhage (the most significant complication, occurring in approximately two to four percent of cases), dehydration from pain-related reduced oral intake, and velopharyngeal insufficiency. Postoperative pain is typically moderate to severe and lasts seven to ten days, and patients should avoid aspirin and non-steroidal anti-inflammatory drugs due to the increased risk of bleeding. The intracapsular technique preserves a thin rim of tonsillar tissue, reducing postoperative pain and hemorrhage rates, though with a small risk of tonsillar regrowth.

\medterm{Tonsillitis}-- Inflammation of the tonsils, most commonly the palatine tonsils, which are lymphoid tissue structures located in the oropharynx between the palatoglossal and palatopharyngeal

\medterm{Turbinates}-- The nasal turbinates, also known as conchae, are three pairs of scroll-like bony projections on the lateral wall of the nasal cavity that are covered by a highly vascularized respiratory mucosa, and they function to increase the surface area available for air conditioning, create turbulent airflow for optimal filtration and warming of inspired air, and regulate nasal airflow through alternating congestion and decongestion of the venous sinusoids within the mucosa (nasal cycle). The inferior turbinate is the largest and is independently mobile, while the middle and superior turbinates are parts of the ethmoid bone. The inferior turbinate mucosa contains large venous sinusoids that can rapidly engorge with blood, causing nasal congestion, and its size fluctuates throughout the day in response to autonomic nervous system input, environmental temperature, allergens, and infection. Chronic enlargement of the inferior turbinates, called turbinate hypertrophy, is a common cause of nasal obstruction and may require medical treatment with topical corticosteroids and antihistamines, or surgical reduction procedures including turbinate cauterization, submucosal resection, or outfracture.

\medterm{Turbinates Hypertrophy}-- Chronic enlargement of the nasal turbinates, most commonly the inferior turbinates, resulting in nasal airway obstruction, which can be caused by chronic inflammation from allergic rhinitis, non-allergic rhinitis, chronic sinusitis, hormonal changes, or vasomotor dysfunction. The inferior turbinates contain large venous sinusoids that can rapidly engorge with blood, and chronic stimulation from allergens, irritants, or infection leads to sustained engorgement and eventual hypertrophy of both the mucosal and bony components of the turbinate. The condition contributes to nasal obstruction, mouth breathing, snoring, and impaired quality of life, and it may be difficult to distinguish from the reversible congestion of acute rhinitis. Medical management includes intranasal corticosteroids, intranasal antihistamines, and avoidance of triggers, while surgical reduction may be considered for refractory cases and includes turbinate outfracture, submucosal resection (using radiofrequency, electrocautery, or microdebrider), and partial turbinectomy, with care taken to avoid excessive turbinate resection which can cause empty nose syndrome.

\medterm{Tympanic membrane}-- The tympanic membrane (eardrum) is a thin, cone-shaped, translucent membrane at the medial end of the external auditory canal that separates the external ear from the middle ear and serves as the interface for sound transmission by converting acoustic energy into mechanical vibrations of the ossicular chain. It is composed of three layers: an outer squamous epithelial layer continuous with the skin of the ear canal, a middle fibrous layer (lamina propria) containing radial and circular collagen fibers that provide structural integrity, and an inner mucosal layer continuous with the lining of the middle ear cavity. The pars tensa constitutes the majority of the membrane and contains a fibrous layer, while the smaller pars flaccida (Shrapnell membrane) lacks a fibrous layer and is more susceptible to retraction pocket formation and cholesteatoma development. Landmarks visible on otoscopy include the handle of the malleus (manubrium), the umbo at its tip, the cone of light reflex in the anteroinferior quadrant, and the pars flaccida superiorly. Perforation of the tympanic membrane may result from trauma, infection, or barotrauma and typically heals spontaneously within weeks, though persistent perforations may require tympanoplasty.

\medterm{Tympanometry}-- An objective, quantitative test that measures the compliance (mobility) of the tympanic membrane and the pressure within the middle ear as a function of air pressure in the external auditory canal, providing information about middle ear function that supplements otoscopic examination. The test uses a tympanometer that delivers a probe tone (typically 226 Hz) into the sealed ear canal while varying the air pressure from positive to negative, measuring the reflected sound energy at each pressure level to generate a tympanogram.

\medterm{Uvula}-- A small, fleshy flap of tissue that hangs from the back of the throat over the root of the tongue.

\medterm{Vertigo}-- A false sensation of rotational or linear movement of the self or the environment, and it is a symptom rather than a disease, requiring differentiation from dizziness, disequilibrium, and presyncope. Peripheral vertigo, caused by dysfunction of the vestibular apparatus or vestibular nerve, is typically characterized by episodic, rotational vertigo that is severe, lasting seconds to hours, often with nausea and vomiting, and associated with nystagmus that is suppressed by visual fixation, hearing loss, and tinnitus. Central vertigo, caused by dysfunction

\medterm{Videostroboscopy}-- A diagnostic technique used to evaluate the structure and function of the vocal cords by producing a slow-motion, composite image of vocal cord vibration using stroboscopic light flashes synchronized to the frequency of vocal cord vibration, allowing the assessment of vocal cord mucosal wave, amplitude, symmetry, phase closure, and

\medterm{Vocal Cord Nodules}-- Vocal cord nodules, also known as singer’s nodules or contact nodules, are bilateral, symmetric, benign epithelial thickenings that develop on the anterior third of the vocal cords at the junction of the membranous and cartilaginous portions, caused by chronic vocal abuse, misuse, or overuse that results in repeated mechanical trauma to the vocal cord mucosa. The nodules are composed of fibrosis, edema, and thickened epithelium in the superficial lamina propria, and they occur at the point of maximum vibration and contact stress on the vocal cords. Patients typically present with hoarseness, breathiness, and vocal fatigue, and the diagnosis is confirmed by laryngoscopy, which reveals bilateral, symmetric, sessile or pedunculated masses on the vocal cords that may touch when the cords are adducted (kissing nodules). Treatment involves voice therapy with modification of vocal habits, avoidance of whispering and shouting, and hydration, and surgical removal is reserved for nodules that are refractory to voice therapy, are large enough to obstruct the glottis, or are suspected of harboring dysplasia.

\medterm{Vocal Cord Paralysis}-- Immobility of one or both vocal cords due to disruption of the recurrent laryngeal nerve (RLN). Unilateral vocal cord paralysis: hoarseness, breathy voice, aspiration, and ineffective cough. Causes: iatrogenic (thyroid, cardiac, spinal, thoracic surgery — most common), malignancy (lung, thyroid, esophageal, mediastinal), trauma, and idiopathic (10-20%). Bilateral vocal cord paralysis: stridor, dyspnea, airway obstruction, but the voice may be normal. Diagnosis: laryngoscopy, CT neck/chest to identify underlying cause, laryngeal EMG. Treatment: unilateral: voice therapy, vocal cord injection (temporary), medialization thyroplasty (permanent). Bilateral: tracheostomy, partial arytenoidectomy, or posterior cordotomy for airway.

\medterm{Vocal Cord Paresis}-- Partial paralysis or weakness of one or both vocal cords due to impaired function of the recurrent laryngeal nerve (RLN) or the vagus nerve. Unilateral paresis is more common and presents with hoarseness, breathy dysphonia, vocal fatigue, and sometimes aspiration. Causes include iatrogenic injury (thyroid, parathyroid, carotid, or thoracic surgery), neoplasms (lung, thyroid, oesophageal), and idiopathic origin. Bilateral paresis presents with inspiratory stridor and dyspnea due to paramedian cord position, often requiring tracheostomy or cord lateralization. Diagnosis: laryngoscopy shows reduced or absent vocal cord movement on the affected side. Laryngeal electromyography can help differentiate neurogenic from mechanical causes. Management depends on severity: voice therapy for mild cases, medialization thyroplasty or injection laryngoplasty for glottic insufficiency, and arytenoid adduction for symptomatic unilateral paresis. Bilateral cases may require posterior cordotomy, arytenoidectomy, or tracheostomy.

\medterm{Vocal Cord Polyp}-- Vocal cord polyps are benign, unilateral, sessile or pedunculated lesions that arise from the vocal cord mucosa, typically as a result of a single traumatic vocal event such as shouting, screaming, or coughing, causing hemorrhage and edema in the superficial lamina propria (Reinke space). Unlike vocal cord nodules, which are bilateral and symmetric, polyps are typically unilateral and may be located anywhere along the vocal cord, though they are most common on the anterior third. The polyps may be translucent (edematous), hemorrhagic, or fibrotic, and they interfere with normal vocal cord vibration, causing hoarseness, breathiness, and a rough voice quality. Diagnosis is confirmed by laryngoscopy, which reveals a smooth, pedunculated or sessile mass protruding from the vocal cord. Treatment involves voice therapy for small polyps, surgical removal for larger polyps or those that are refractory to voice therapy, and avoidance of the vocal abuse that caused the initial trauma to prevent recurrence.

\medterm{Vocal Cords}-- The vocal cords, also known as vocal folds, are two pairs of horizontal folds of tissue within the

\medterm{Voice and Speech}-- The products of complex coordinated activity of the respiratory, laryngeal, and articulatory systems. Voice (phonation) is produced by airflow from the lungs passing through the larynx causing vibration of the vocal cords (vocal folds); fundamental frequency (pitch) is determined by vocal cord length and tension, intensity (loudness) by subglottic air pressure, and quality by the pattern of vocal cord vibration. Speech (articulation) involves modification of the laryngeal sound by the pharynx, oral cavity, tongue, lips, teeth, and palate (resonance). Voice disorders (dysphonia): organic (vocal cord nodules, polyps, Reinke oedema, laryngitis, laryngeal cancer, vocal cord paralysis/paresis) or functional (muscle tension dysphonia, psychogenic dysphonia, puberphonia). Speech disorders: dysarthria (neurological weakness/incoordination of articulators — stroke, Parkinson disease, ALS, cerebellar disease), apraxia of speech (motor planning deficit), and articulation disorders. Voice/speech assessment: laryngoscopy (direct or videostroboscopy) and acoustic analysis.

\medterm{Weber Test}-- The Weber test is a screening test for lateralization of hearing that uses a vibrating tuning fork (typically 512 Hz) placed on the midline of the skull (vertex or forehead) to determine whether the sound is heard equally in both ears or louder in one ear, providing information about the type and side of hearing loss. In a patient with normal hearing, the sound is perceived equally in both ears or slightly louder in one ear due to normal asymmetry. In conductive hearing loss, the sound lateralizes to the affected ear because ambient noise is reduced on that side, enhancing the perception of bone-conducted sound. In sensorineural hearing loss, the sound lateralizes to the unaffected ear because the damaged cochlea is less sensitive to the bone-conducted sound.

\medterm{Xerostomia (Dry Mouth)}-- A subjective sensation of oral dryness resulting from reduced or absent saliva production (salivary gland hypofunction). Xerostomia affects approximately 20% of the general population, increasing with age and with the number of medications taken. Causes: medications (the most common cause, over 500 drugs cause xerostomia — anticholinergics, tricyclic antidepressants, SSRIs, antipsychotics, antihistamines, decongestants, antihypertensives, diuretics, opioids, and proton pump inhibitors), radiation therapy (head and neck cancer radiotherapy causes irreversible damage to salivary glands, dose-dependent, >30 Gy causes permanent hypofunction), systemic diseases (Sjogren syndrome, an autoimmune exocrinopathy affecting lacrimal and salivary glands; diabetes mellitus; HIV/AIDS; sarcoidosis; amyloidosis), and dehydration (poor fluid intake, vomiting, diarrhoea, polyuria). Consequences: increased risk of dental caries (root caries, smooth surface caries), oral candidiasis (thrush), dysphagia (difficulty swallowing dry foods), dysgeusia (taste disturbance), oral mucosal discomfort and burning, halitosis, and difficulty wearing dentures. Diagnosis: sialometry (unstimulated whole salivary flow rate <0.1 mL/min is hyposalivation), minor salivary gland biopsy for suspected Sjogren syndrome. Treatment: symptomatic — frequent water sips, sugar-free chewing gum (stimulates residual salivary flow), artificial saliva substitutes (carboxymethylcellulose, mucin-based sprays, gels, and lozenges), and salivary stimulants (pilocarpine 5 mg tid, cevimeline 30 mg tid — muscarinic agonists that stimulate residual salivary tissue, contraindicated in uncontrolled asthma, narrow-angle glaucoma, and iritis). Preventive: meticulous oral hygiene, regular fluoride application (high-fluoride toothpaste, fluoride varnish), chlorhexidine mouthwash for caries prevention, and antifungal therapy for oral candidiasis. that uses a vibrating tuning fork (typically 512 Hz) placed on the midline of the skull (vertex or forehead) to determine whether the sound is heard equally in both ears or louder in one ear, providing information about the type and side of hearing loss. In a patient with normal hearing, the sound is perceived equally in both ears or slightly louder in one ear due to normal asymmetry. In conductive hearing loss, the sound lateralizes to the affected ear because ambient noise is reduced on that side, enhancing the perception of bone-conducted sound. In sensorineural hearing loss, the sound lateralizes to the unaffected ear because the damaged cochlea is less sensitive to the bone-conducted sound.
